\section{Experiment}
\label{sec:experiment}

\subsection{Evaluation Metrics}
\label{sec:metrics}

To ensure a rigorous validation of the proposed framework, particularly regarding its clinical utility and robustness against class imbalance, we employ a comprehensive suite of evaluation metrics. These metrics encompass classification performance, ranking quality, inter-rater reliability, and model calibration.

\paragraph{Classification Metrics}
Given the class imbalance in the dataset, we adopt complementary metrics to provide a holistic view of model performance.

\vspace{0.5em}

The Top-1 Accuracy measures overall classification correctness:

\begin{equation}
\text{Accuracy} = \frac{1}{N} \sum_{i=1}^{N} \mathbb{I}[\hat{y}_i = y_i]
\end{equation}

where $N$ is the number of test samples and $\mathbb{I}[\cdot]$ is the indicator function.

The Macro-averaged F1-score (F1-macro) computes the F1-score independently for each class and then calculates the unweighted mean. This metric is critical for our study as it treats all disease categories equally, ensuring that performance on minority conditions (e.g., Hypodontia) is not overshadowed by majority classes~\cite{tsoumakas2007multi}:

\begin{equation}
\text{F1-macro} = \frac{1}{Q} \sum_{j=1}^{Q} \frac{2 TP_j}{2 TP_j + FP_j + FN_j}
\end{equation}

where $Q = 6$ denotes the number of disease classes.

\paragraph{Ranking Quality}
To evaluate the quality of predicted probability rankings across classes, we employ mean Average Precision (mAP). For multi-class classification, mAP measures how well the model ranks the correct class among alternatives~\cite{zhang2014review}:

\begin{equation}
\text{mAP} = \frac{1}{Q} \sum_{j=1}^{Q} \text{AP}_j, \quad \text{AP}_j = \int_{0}^{1} p_j(r) \, dr
\end{equation}

where $p_j(r)$ is the precision-recall curve for class $j$. mAP is threshold-independent and provides insights into model confidence calibration across all disease categories.

\paragraph{Reliability and Calibration}
To assess the agreement between model predictions and ground truth beyond random chance, we employ Cohen's Kappa ($\kappa$). For the multi-class classification setting, we compute the macro-averaged Kappa by treating each class as a binary classification problem and reporting the mean across classes~\cite{powers2011evaluation}:

\begin{equation}
\kappa_j = \frac{p_{o,j} - p_{e,j}}{1 - p_{e,j}}, \quad \kappa_{\text{macro}} = \frac{1}{Q} \sum_{j=1}^{Q} \kappa_j
\end{equation}

where $p_{o,j}$ represents observed agreement and $p_{e,j}$ represents expected chance agreement. A $\kappa_{\text{macro}} > 0.80$ is typically required for reliable clinical deployment.

Additionally, model trustworthiness is evaluated using the Expected Calibration Error (ECE). ECE measures the weighted average difference between the model's confidence and its actual accuracy across $M$ bins (set to $M=15$)~\cite{naeini2015obtaining}:

\begin{equation}
\text{ECE} = \sum_{m=1}^{M} \frac{|B_m|}{N} |\text{acc}(B_m) - \text{conf}(B_m)|
\end{equation}

Low ECE values ($< 0.05$) indicate that the model's probability estimates accurately reflect its true predictive correctness.

\paragraph{Robustness Assessment}
Finally, to simulate real-world deployment challenges, we evaluate performance stability under image perturbations (e.g., Gaussian noise, lighting variations, motion blur, JPEG compression artifacts). We define the Robustness Retention Rate (RR) for a given metric $\mathcal{M}$ as~\cite{hendrycks2019benchmarking}:

\begin{equation}
\text{RR}_{\mathcal{M}} = \frac{\mathcal{M}_{\text{perturbed}}}{\mathcal{M}_{\text{clean}}} \times 100\%
\end{equation}

We specifically report $\text{RR}_{\text{F1-macro}}$ and $\text{RR}_{\text{Accuracy}}$ to ensure the classification system remains reliable under adverse imaging conditions. Models achieving $\text{RR} > 90\%$ under moderate perturbations are considered deployment-ready.

\subsection{Ablation Study}
\label{sec:ablation}

To validate the contribution of each component and justify key hyperparameter choices, we conduct systematic ablation experiments. All experiments use identical training configurations unless otherwise noted.

\subsubsection{Hyperparameter Sensitivity Analysis}

We evaluate the sensitivity of three critical hyperparameters: window size $M$ for self-attention, alpha initialization $\alpha$ for gated fusion, and the learning rate ratio between backbone and classification head.

\begin{table}[htbp]
\centering
\caption{Hyperparameter Sensitivity Analysis. Best values in \textbf{bold}.}
\label{tab:hyperparams}
\small
\resizebox{\columnwidth}{!}{%
\begin{tabular}{llccc}
\toprule
\textbf{Param} & \textbf{Value} & \textbf{F1} & \textbf{Acc} & \textbf{Notes} \\
\midrule
\multirow{3}{*}{Window $M$} 
& 5 & XX.X\% & XX.X\% & 25 tok/win \\
& \textbf{7} & \textbf{XX.X\%} & \textbf{XX.X\%} & 49 tok/win \\
& 9 & XX.X\% & XX.X\% & 81 tok/win \\
\midrule
\multirow{4}{*}{Alpha $\alpha$}
& 0.05 & XX.X\% & XX.X\% & Conservative \\
& \textbf{0.1} & \textbf{XX.X\%} & \textbf{XX.X\%} & Balanced \\
& 0.2 & XX.X\% & XX.X\% & Aggressive \\
& 0.5 & XX.X\% & XX.X\% & Equal weight \\
\midrule
\multirow{3}{*}{LR Ratio}
& 1:10 & XX.X\% & XX.X\% & Fast backbone \\
& 1:20 & XX.X\% & XX.X\% & Moderate \\
& \textbf{1:30} & \textbf{XX.X\%} & \textbf{XX.X\%} & Preserve pretrained \\
\bottomrule
\end{tabular}%
}
\end{table}

\textbf{Window Size $M$}: We test $M \in \{5, 7, 9\}$ (Table~\ref{tab:hyperparams}). $M=7$ achieves optimal performance, balancing local context (49 tokens per window) with computational efficiency. Smaller windows ($M=5$) limit receptive field, while larger windows ($M=9$) increase FLOPs without proportional accuracy gains.

\textbf{Alpha Initialization}: The gated fusion parameter $\alpha$ controls CNN-Transformer feature balance (Eq.~\ref{eq:fusion}). We initialize $\alpha=0.1$ following~\cite{touvron2021training}, which prevents gradient instability in early training while allowing gradual feature fusion. Higher initial values ($\alpha=0.5$) cause training oscillation; lower values ($\alpha=0.05$) slow convergence.

\textbf{Learning Rate Ratio}: Differential learning rates preserve pretrained backbone features while allowing the classification head to adapt quickly. LR ratio 1:30 (backbone $1\times10^{-5}$, head $3\times10^{-4}$) optimally balances feature preservation with domain adaptation.

\subsubsection{Backbone Architecture Comparison}

\begin{table}[htbp]
\centering
\caption{Backbone Architecture Comparison}
\label{tab:ablation_backbone}
\small
\resizebox{\columnwidth}{!}{%
\begin{tabular}{lcccc}
\toprule
\textbf{Backbone} & \textbf{Params} & \textbf{F1} & \textbf{Acc} & \textbf{Inf. (ms)} \\
\midrule
MobileNetV3-S & 8.2M & XX.X\% & XX.X\% & 28 \\
\textbf{EffNet-B0} & \textbf{15.7M} & \textbf{XX.X\%} & \textbf{XX.X\%} & \textbf{42} \\
EffNet-B3 & 27.4M & XX.X\% & XX.X\% & 89 \\
ResNet-50 & 31.2M & XX.X\% & XX.X\% & 67 \\
\bottomrule
\end{tabular}%
}
\end{table}

EfficientNet-B0 achieves optimal accuracy-efficiency trade-off (Table~\ref{tab:ablation_backbone}). While EfficientNet-B3 provides marginally higher accuracy, the 2$\times$ inference time increase is impractical for clinical deployment.

\subsubsection{Loss Function Comparison}

\begin{table}[htbp]
\centering
\caption{Loss Function Comparison}
\label{tab:ablation_loss}
\small
\resizebox{\columnwidth}{!}{%
\begin{tabular}{lccc}
\toprule
\textbf{Loss Function} & \textbf{F1-macro} & \textbf{F1-min} & \textbf{ECE} \\
\midrule
Cross-Entropy & XX.X\% & XX.X\% & 0.XXX \\
Focal ($\gamma$=2) & XX.X\% & XX.X\% & 0.XXX \\
Label Smooth ($\epsilon$=0.1) & XX.X\% & XX.X\% & 0.XXX \\
\textbf{Focal + LS} & \textbf{XX.X\%} & \textbf{XX.X\%} & \textbf{0.XXX} \\
\bottomrule
\end{tabular}%
}
\end{table}

Focal Loss with label smoothing achieves the best F1-macro and calibration (Table~\ref{tab:ablation_loss}). F1-minority reports average F1 on minority classes (Calculus, Hypodontia), showing particular improvement with focal loss.

\subsubsection{Component Contribution Analysis}

\begin{table}[htbp]
\centering
\caption{Component Contribution Analysis (Progressive Ablation)}
\label{tab:ablation_components}
\small
\resizebox{\columnwidth}{!}{%
\begin{tabular}{lccc}
\toprule
\textbf{Configuration} & \textbf{F1-macro} & \textbf{$\Delta$ F1} & \textbf{Params} \\
\midrule
CNN Only (EffNet-B0) & XX.X\% & -- & 5.3M \\
+ Transformer (parallel) & XX.X\% & +X.X\% & 12.1M \\
+ Additive Fusion & XX.X\% & +X.X\% & 14.3M \\
+ Gated Cross-Attn (Ours) & \textbf{XX.X\%} & +X.X\% & 15.7M \\
\bottomrule
\end{tabular}%
}
\end{table}

Progressive ablation (Table~\ref{tab:ablation_components}) demonstrates that each component contributes to overall performance. The Transformer branch adds global context modeling, additive fusion enables feature interaction, and gated cross-attention provides adaptive local-global balance.